% Ban thu: chi giu cac cau co trong CD2

%@ id=TT25.14.003 | trang=1 | so=2 | de=II | vitri=CD2.4,C3.9.5 | muc=8 | tt=0 | soat=ok
\begin{bt}
Với $x > 0$, cho hai biểu thức $A = \dfrac{x + 2\sqrt{x}}{x}$ và $B = \dfrac{\sqrt{x}-1}{\sqrt{x}} + \dfrac{2\sqrt{x}+1}{x+\sqrt{x}}$.
\begin{cacy}
\item Tính giá trị của biểu thức $A$ khi $x = 4$.
\item Chứng minh $B = \dfrac{\sqrt{x}+2}{\sqrt{x}+1}$.
\item Tìm số nguyên $x$ nhỏ nhất để $\dfrac{A}{B} < \dfrac{7}{4}$.
\end{cacy}
\loigiai{
\begin{cacy}
\item Thay $x = 4$ (thỏa mãn điều kiện) vào biểu thức $A$, ta được:

Ta có $A = \dfrac{4 + 2\sqrt{4}}{4} = 2$.

Vậy $A = 2$ khi $x = 4$.

\item Ta có:
\[
B = \dfrac{\sqrt{x}-1}{\sqrt{x}} + \dfrac{2\sqrt{x}+1}{x+\sqrt{x}} = \dfrac{(\sqrt{x}-1)(\sqrt{x}+1) + 2\sqrt{x}+1}{\sqrt{x}(\sqrt{x}+1)} = \dfrac{x+2\sqrt{x}}{\sqrt{x}(\sqrt{x}+1)} = \dfrac{\sqrt{x}+2}{\sqrt{x}+1}.
\]
Vậy với $x > 0$ thì $B = \dfrac{\sqrt{x}+2}{\sqrt{x}+1}$.

\item Với $x > 0$,
\[
\dfrac{A}{B} = \dfrac{x+2\sqrt{x}}{x} : \dfrac{\sqrt{x}+2}{\sqrt{x}+1} = \dfrac{\sqrt{x}(\sqrt{x}+2)}{x} \cdot \dfrac{\sqrt{x}+1}{\sqrt{x}+2} = \dfrac{\sqrt{x}(\sqrt{x}+1)}{x} = \dfrac{x+\sqrt{x}}{x}.
\]
Ta có:
\[ \dfrac{A}{B} < \dfrac{7}{4} \]
\[ \dfrac{x+\sqrt{x}}{x} < \dfrac{7}{4} \]
\[ 1 + \dfrac{1}{\sqrt{x}} < \dfrac{7}{4} \]
\[ \dfrac{1}{\sqrt{x}} < \dfrac{3}{4} \]
\[ \sqrt{x} > \dfrac{4}{3} \]
\[ x > \dfrac{16}{9}. \]
Kết hợp điều kiện $x > 0$ và $x$ là số nguyên nhỏ nhất nên $x = 2$.

Vậy $x = 2$.
\end{cacy}
}
\end{bt}
